% ============================================================
% CHAPITRE 12 — Texte et effets graphiques dans Photoshop
% ============================================================

\chapter{Texte et effets graphiques dans Photoshop}

\noindent
\textit{Dans les chapitres précédents, nous avons appris à créer des documents, à organiser des calques, à sélectionner, à détourer, à retoucher des photographies, à réaliser des photomontages. Toutes ces compétences sont essentielles pour composer des visuels. Mais un visuel sans texte est souvent incomplet : c'est le texte qui porte le message, identifie la marque, donne les informations. Il faut donc apprendre à intégrer du texte dans Photoshop, et à lui appliquer des effets graphiques.}

\vspace{1em}

\noindent
Le \terme{texte} est l'un des éléments centraux du design graphique. Il porte le message verbal, structure la composition, transmet une identité. Dans Photoshop, le texte est un calque spécifique, modifiable tant qu'il n'est pas vectorisé. On peut lui appliquer des propriétés typographiques (police, taille, graisse, couleur, alignement), mais aussi des effets graphiques (ombres, contours, dégradés, textures, profondeur). Ces effets permettent de créer des visuels typographiques spectaculaires, utilisés dans la publicité, l'édition, le web, le packaging.

\vspace{1em}

\noindent
Mais les effets graphiques sont à double tranchant : bien utilisés, ils subliment le texte ; mal utilisés, ils le rendent illisible, kitsch, amateur. Il faut donc apprendre à les utiliser avec \terme{maîtrise} : un effet doit servir le message, pas l'inverse. Un texte avec trop d'effets perd sa lisibilité et sa crédibilité.

\vspace{1em}

\noindent
Ce chapitre a pour objectif de te faire maîtriser le texte et les effets graphiques dans Photoshop. Nous commencerons par apprendre à créer du texte. Nous verrons ensuite comment modifier les propriétés typographiques, organiser une hiérarchie de titres, utiliser les styles de calque, créer des ombres portées, appliquer des contours et des dégradés, créer des effets de profondeur, intégrer des textures, construire une composition publicitaire, et enfin éviter la surcharge d'effets. À la fin de ce chapitre, tu sauras créer des visuels typographiques professionnels.

% ============================================================
\section{Créer du texte}
% ============================================================

\subsection{Pourquoi créer du texte dans Photoshop ?}

Photoshop n'est pas le logiciel idéal pour la typographie pure (InDesign est plus adapté pour la mise en page éditoriale). Mais Photoshop est très utilisé pour :

\begin{itemize}
    \item \terme{Créer} des visuels publicitaires (affiches, bannières, posts réseaux sociaux).
    \item \terme{Intégrer} du texte dans une composition graphique (photomontage, collage).
    \item \terme{Appliquer} des effets graphiques au texte (ombres, contours, textures).
    \item \terme{Créer} des titres, des slogans, des accroches.
    \item \terme{Préparer} des maquettes de sites web, d'applications.
\end{itemize}

Photoshop propose des outils de texte complets, qui permettent de créer, modifier, styliser du texte de manière professionnelle.

\subsection{Les outils de texte}

\subsubsection{L'outil Texte horizontal}

L'\terme{outil Texte horizontal} (\texttt{T}) permet de créer du texte horizontal. C'est l'outil le plus utilisé.

\subsubsection{L'outil Texte vertical}

L'\terme{outil Texte vertical} permet de créer du texte vertical (caractères empilés verticalement). Il est utilisé pour des usages spécifiques (affiches japonaises, titres verticaux).

\subsubsection{L'outil Texte masque horizontal}

L'\terme{outil Texte masque horizontal} permet de créer une sélection en forme de texte. Utile pour appliquer des effets à un texte sans créer de calque de texte.

\subsubsection{L'outil Texte masque vertical}

L'\terme{outil Texte masque vertical} permet de créer une sélection en forme de texte vertical.

\subsection{Les types de texte}

\subsubsection{Le texte ponctuel}

Le \terme{texte ponctuel} est un texte créé par un simple clic sur la zone de travail. Il n'a pas de zone délimitée : il s'étend sur une ligne (ou plusieurs lignes si on appuie sur \texttt{Entrée}). Il est utilisé pour :

\begin{itemize}
    \item Les titres.
    \item Les slogans.
    \item Les mots isolés.
    \item Les étiquettes.
\end{itemize}

\subsubsection{Le texte de paragraphe}

Le \terme{texte de paragraphe} est un texte créé en cliquant-glissant pour définir une zone rectangulaire. Il s'étend dans cette zone, avec des retours à la ligne automatiques. Il est utilisé pour :

\begin{itemize}
    \item Les paragraphes de texte.
    \item Les descriptions.
    \item Les mentions légales.
    \item Les blocs de contenu.
\end{itemize}

\subsubsection{Le texte curviligne}

Le \terme{texte curviligne} est un texte créé sur un tracé (courbe, forme). Il suit la courbe du tracé. Il est utilisé pour :

\begin{itemize}
    \item Les logos circulaires.
    \item Les badges.
    \item Les tampons.
    \item Les compositions créatives.
\end{itemize}

\subsection{Créer un texte ponctuel}

\subsubsection{Étapes}

\begin{enumerate}
    \item Sélectionner l'outil Texte horizontal (\texttt{T}).
    \item Cliquer sur la zone de travail.
    \item Saisir le texte.
    \item Valider avec \texttt{Ctrl + Entrée} (ou \texttt{Cmd + Entrée}), ou cliquer sur la coche dans la barre d'options.
\end{enumerate}

\subsubsection{Le calque de texte}

Un calque de texte est automatiquement créé. Il est reconnaissable à l'icône \textbf{T} dans le panneau Calques. Le texte est modifiable tant qu'il n'est pas vectorisé.

\subsection{Créer un texte de paragraphe}

\subsubsection{Étapes}

\begin{enumerate}
    \item Sélectionner l'outil Texte horizontal (\texttt{T}).
    \item Cliquer-glisser pour définir une zone rectangulaire.
    \item Saisir le texte.
    \item Valider avec \texttt{Ctrl + Entrée} (ou \texttt{Cmd + Entrée}).
\end{enumerate}

\subsubsection{Le cadre de texte}

Le texte de paragraphe est contenu dans un cadre. On peut :

\begin{itemize}
    \item \terme{Redimensionner} le cadre en tirant sur les poignées.
    \item \terme{Déplacer} le cadre en cliquant-glissant à l'intérieur.
    \item \terme{Dépasser} le cadre : si le texte est trop long, un signe \textbf{+} apparaît en bas à droite.
\end{itemize}

\subsection{Créer un texte curviligne}

\subsubsection{Étapes}

\begin{enumerate}
    \item Créer un tracé (avec l'outil Plume, ou une forme).
    \item Sélectionner l'outil Texte horizontal (\texttt{T}).
    \item Cliquer sur le tracé.
    \item Saisir le texte.
    \item Valider.
\end{enumerate}

\subsubsection{Le tracé de texte}

Le texte suit le tracé. On peut :

\begin{itemize}
    \item \terme{Déplacer} le texte le long du tracé.
    \item \terme{Retourner} le texte (de l'autre côté du tracé).
    \item \terme{Modifier} le tracé pour changer la courbe.
\end{itemize}

\subsection{Modifier un texte}

\subsubsection{Modifier le contenu}

\begin{enumerate}
    \item Sélectionner l'outil Texte (\texttt{T}).
    \item Cliquer sur le texte existant.
    \item Modifier le contenu.
    \item Valider.
\end{enumerate}

\subsubsection{Modifier les propriétés}

\begin{itemize}
    \item \terme{Police} : changer la police.
    \item \terme{Taille} : changer la taille.
    \item \terme{Couleur} : changer la couleur.
    \item \terme{Graisse} : changer la graisse (normal, gras, etc.).
    \item \terme{Style} : changer le style (italique, souligné).
    \item \terme{Alignement} : changer l'alignement (gauche, centre, droite, justifié).
\end{itemize}

\subsection{Vectoriser un texte}

\subsubsection{Principe}

\terme{Vectoriser} un texte consiste à le transformer en forme vectorielle. Le texte n'est plus modifiable, mais il peut être déformé, stylisé, transformé comme une forme.

\subsubsection{Utilisation}

Clic droit sur le calque de texte > \texttt{Convertir en forme}, ou \texttt{Calque > Vectoriser le texte}.

\subsubsection{Précautions}

\begin{itemize}
    \item \terme{Conserver} une copie du calque de texte avant de vectoriser.
    \item \terme{Vectoriser} uniquement si nécessaire.
    \item \terme{Ne pas vectoriser} un texte destiné à être modifié.
\end{itemize}

\subsection{Les bonnes pratiques}

\begin{itemize}
    \item \terme{Utiliser} l'outil Texte pour créer du texte.
    \item \terme{Choisir} le type de texte adapté (ponctuel, paragraphe, curviligne).
    \item \terme{Nommer} le calque de texte.
    \item \terme{Conserver} le texte modifiable (ne pas vectoriser sans nécessité).
    \item \terme{Vérifier} la lisibilité (taille, contraste).
    \item \terme{Utiliser} les polices adaptées.
    \item \terme{Enregistrer} régulièrement.
\end{itemize}

\subsection*{Conclusion de la section}

Créer du texte dans Photoshop est une opération courante. On utilise l'outil Texte (\texttt{T}) pour créer du texte ponctuel, du texte de paragraphe, ou du texte curviligne. Le texte est un calque spécifique, modifiable tant qu'il n'est pas vectorisé. On peut modifier le contenu, les propriétés, l'alignement. Il est essentiel de nommer les calques, de conserver le texte modifiable, de vérifier la lisibilité.

Maintenant que nous savons créer du texte, nous allons, dans la section suivante, apprendre à \terme{modifier les propriétés typographiques}.

% ============================================================
\section{Modifier les propriétés typographiques}
% ============================================================

\subsection{Les panneaux de typographie}

\subsubsection{Le panneau Caractère}

Le panneau \terme{Caractère} permet de modifier les propriétés typographiques du texte sélectionné :

\begin{itemize}
    \item \terme{Police} : choisir la police.
    \item \terme{Style} : normal, gras, italique, etc.
    \item \terme{Taille} : taille du texte.
    \item \terme{Interlignage} : espacement entre les lignes.
    \item \terme{Crénage} : espacement entre les caractères.
    \item \terme{Approche} : espacement global entre les caractères.
    \item \terme{Couleur} : couleur du texte.
    \item \terme{Majuscules/Minuscules} : petites majuscules, exposant, indice.
    \item \terme{Soulignement, barré} : styles de décoration.
    \item \terme{Langue} : langue du texte (pour la césure et la correction).
\end{itemize}

\subsubsection{Le panneau Paragraphe}

Le panneau \terme{Paragraphe} permet de modifier les propriétés de mise en page du texte :

\begin{itemize}
    \item \terme{Alignement} : gauche, centre, droite, justifié.
    \item \terme{Retrait} : retrait de première ligne, retrait à gauche, retrait à droite.
    \item \terme{Espacement} : espacement avant et après les paragraphes.
    \item \terme{Césure} : césure automatique.
    \item \terme{Ligne de base} : décalage de la ligne de base.
    \item \terme{Composition} : options de composition.
\end{itemize}

\subsection{La barre d'options}

\subsubsection{Les options de base}

Lorsque l'outil Texte est sélectionné, la barre d'options propose :

\begin{itemize}
    \item \terme{Police} : choisir la police.
    \item \terme{Style} : normal, gras, italique.
    \item \terme{Taille} : taille du texte.
    \item \terme{Alignement} : gauche, centre, droite, justifié.
    \item \terme{Couleur} : couleur du texte.
    \item \terme{Déformation} : déformer le texte.
    \item \terme{Panoramique} : créer un texte.
    \item \terme{Paragraphe} : créer un texte de paragraphe.
\end{itemize}

\subsubsection{Les options avancées}

\begin{itemize}
    \item \terme{Anti-aliasing} : lissage des bords.
    \item \terme{Résolution} : résolution du texte.
    \item \terme{Alignement vertical} : haut, centre, bas.
\end{itemize}

\subsection{La couleur du texte}

\subsubsection{Changer la couleur}

\begin{enumerate}
    \item Sélectionner le texte.
    \item Cliquer sur la couleur dans la barre d'options ou le panneau Caractère.
    \item Choisir une couleur dans le sélecteur.
    \item Valider.
\end{enumerate}

\subsubsection{Appliquer une couleur différente à une partie du texte}

\begin{enumerate}
    \item Sélectionner une partie du texte avec l'outil Texte.
    \item Changer la couleur.
    \item Valider.
\end{enumerate}

\subsection{L'interlignage, le crénage, l'approche}

\subsubsection{L'interlignage}

L'\terme{interlignage} est l'espace vertical entre les lignes. Il se mesure en points ou en pourcentage de la taille du texte.

\begin{itemize}
    \item \terme{Interlignage automatique} : calculé par Photoshop.
    \item \terme{Interlignage manuel} : défini par l'utilisateur.
    \item \terme{Valeurs courantes} : 120\% à 150\% de la taille du texte.
\end{itemize}

\subsubsection{Le crénage}

Le \terme{crénage} est l'espacement entre deux caractères spécifiques. Il se mesure en millièmes d'em.

\begin{itemize}
    \item \terme{Crénage métrique} : selon les métriques de la police.
    \item \terme{Crénage optique} : ajustement automatique.
    \item \terme{Crénage manuel} : ajustement manuel.
\end{itemize}

\subsubsection{L'approche}

L'\terme{approche} est l'espacement global entre tous les caractères. Elle se mesure en millièmes d'em.

\begin{itemize}
    \item \terme{Approche négative} : les caractères sont rapprochés.
    \item \terme{Approche positive} : les caractères sont espacés.
\end{itemize}

\subsection{Les styles de texte}

\subsubsection{Les styles de paragraphe}

Les \terme{styles de paragraphe} permettent d'enregistrer des propriétés de texte pour les réutiliser. On peut créer un style pour :

\begin{itemize}
    \item Les titres.
    \item Les sous-titres.
    \item Le texte courant.
    \item Les légendes.
\end{itemize}

\subsubsection{Les styles de caractère}

Les \terme{styles de caractère} permettent d'enregistrer des propriétés de caractère (police, taille, couleur) pour les réutiliser.

\subsubsection{Création d'un style}

\begin{enumerate}
    \item Sélectionner le texte avec les propriétés souhaitées.
    \item Ouvrir le panneau \emph{Styles de paragraphe} ou \emph{Styles de caractère}.
    \item Cliquer sur \emph{Créer un style}.
    \item Nommer le style.
    \item Valider.
\end{enumerate}

\subsection{Les bonnes pratiques}

\begin{itemize}
    \item \terme{Choisir} des polices lisibles.
    \item \terme{Limiter} le nombre de polices (1 à 2).
    \item \terme{Utiliser} des tailles cohérentes (échelle typographique).
    \item \terme{Ajuster} l'interlignage pour la lisibilité.
    \item \terme{Utiliser} des styles pour la cohérence.
    \item \terme{Vérifier} la lisibilité sur différents fonds.
    \item \terme{Ne pas abuser} des effets (voir section 10).
\end{itemize}

\subsection*{Conclusion de la section}

Les \terme{propriétés typographiques} permettent de contrôler l'apparence du texte : police, taille, style, interlignage, crénage, approche, couleur, alignement. On les modifie via la barre d'options, les panneaux Caractère et Paragraphe. On peut enregistrer des styles pour la cohérence. Il est essentiel de choisir des polices lisibles, de limiter le nombre de polices, d'ajuster l'interlignage, d'utiliser des styles, de vérifier la lisibilité.

Maintenant que nous savons modifier les propriétés typographiques, nous allons, dans la section suivante, apprendre à \terme{organiser une hiérarchie de titres}.

% ============================================================
\section{Organiser une hiérarchie de titres}
% ============================================================

\subsection{L'importance de la hiérarchie}

La \terme{hiérarchie typographique} organise les différents niveaux de texte selon leur importance :

\begin{itemize}
    \item \terme{Titre principal} : le message le plus important.
    \item \terme{Sous-titre} : les messages secondaires.
    \item \terme{Texte courant} : le contenu.
    \item \terme{Légendes} : les textes accompagnant les images.
    \item \terme{Mentions} : les informations techniques.
\end{itemize}

Une hiérarchie claire guide le regard du lecteur, facilite la compréhension, met en valeur les informations importantes.

\subsection{Les moyens de créer la hiérarchie}

\subsubsection{La taille}

Plus un texte est grand, plus il est important. Le titre principal est le plus grand ; les mentions sont les plus petites.

\subsubsection{La graisse}

Plus un texte est gras, plus il est important. Le titre principal est souvent en gras.

\subsubsection{La couleur}

Les couleurs vives ou contrastées attirent plus que les couleurs neutres. Le titre principal peut être coloré, le texte courant en noir.

\subsubsection{La police}

Changer de police crée une distinction. On peut utiliser une police pour les titres, une autre pour le texte courant.

\subsubsection{La casse}

Les majuscules attirent plus que les minuscules. Le titre peut être en majuscules, le texte courant en minuscules.

\subsubsection{La position}

Les éléments en haut, au centre, à gauche sont vus en premier.

\subsubsection{L'espacement}

Un élément isolé par du vide attire plus.

\subsection{Construire une hiérarchie}

\subsubsection{Étape 1 : Identifier les niveaux}

Déterminer les niveaux de texte nécessaires :

\begin{itemize}
    \item Titre principal.
    \item Sous-titre.
    \item Texte courant.
    \item Légendes.
    \item Mentions.
\end{itemize}

\subsubsection{Étape 2 : Définir les rôles}

Attribuer une police, une taille, une graisse, une couleur à chaque niveau.

\subsubsection{Étape 3 : Créer les styles}

Créer des styles de paragraphe ou de caractère pour chaque niveau.

\subsubsection{Étape 4 : Appliquer les styles}

Appliquer les styles aux textes correspondants.

\subsubsection{Étape 5 : Vérifier}

Vérifier la hiérarchie sur l'ensemble de la composition.

\subsection{Exemple de hiérarchie}

Prenons une affiche de concert :

\begin{itemize}
    \item \terme{Nom du groupe} : 72 pt, gras, noir, en haut.
    \item \terme{Date et lieu} : 24 pt, normal, rouge, au milieu.
    \item \terme{Description} : 14 pt, normal, gris, en bas.
    \item \terme{Mentions légales} : 8 pt, normal, gris clair, tout en bas.
\end{itemize}

Cette hiérarchie guide le regard du nom du groupe vers les mentions légales.

\subsection{Les bonnes pratiques}

\begin{itemize}
    \item \terme{Limiter} le nombre de niveaux (3 à 5).
    \item \terme{Définir} des tailles cohérentes (échelle typographique).
    \item \terme{Utiliser} des polices adaptées à chaque niveau.
    \item \terme{Créer} des styles pour la cohérence.
    \item \terme{Vérifier} la lisibilité.
    \item \terme{Vérifier} la hiérarchie sur différents fonds.
    \item \terme{Ne pas abuser} des contrastes.
\end{itemize}

\subsection*{Conclusion de la section}

La \terme{hiérarchie typographique} organise les niveaux de texte selon leur importance. Elle se crée par la taille, la graisse, la couleur, la police, la casse, la position, l'espacement. On construit une hiérarchie en identifiant les niveaux, définissant les rôles, créant des styles, appliquant les styles, vérifiant. Il est essentiel de limiter le nombre de niveaux, de définir des tailles cohérentes, de créer des styles, de vérifier la lisibilité.

Maintenant que nous savons organiser une hiérarchie, nous allons, dans la section suivante, apprendre à \terme{utiliser les styles de calque}.

% ============================================================
\section{Utiliser les styles de calque}
% ============================================================

\subsection{Qu'est-ce qu'un style de calque ?}

\begin{quote}
\textbf{Un style de calque est un ensemble d'effets appliqués à un calque (ombre, contour, biseau, etc.), qui peuvent être modifiés, désactivés, supprimés à tout moment.}
\end{quote}

Les styles de calque sont un outil puissant et non destructif pour appliquer des effets graphiques au texte (et à d'autres calques). Ils sont très utilisés pour les visuels typographiques.

\subsection{Les avantages des styles de calque}

\begin{itemize}
    \item \terme{Non destructif} : les effets ne modifient pas les pixels du calque.
    \item \terme{Modifiable} : on peut ajuster les paramètres à tout moment.
    \item \terme{Désactivable} : on peut désactiver un effet pour comparer.
    \item \terme{Supprimable} : on peut supprimer un effet.
    \item \terme{Cumulable} : on peut empiler plusieurs effets.
    \item \terme{Copiable} : on peut copier les styles d'un calque à un autre.
\end{itemize}

\subsection{Accéder aux styles de calque}

\subsubsection{Par double-clic}

Double-cliquer sur le calque (à côté du nom) pour ouvrir la boîte de dialogue \emph{Style de calque}.

\subsubsection{Par le bouton fx}

Cliquer sur le bouton \terme{fx} en bas du panneau Calques, puis choisir un effet.

\subsubsection{Par le menu}

\texttt{Calque > Style de calque}, puis choisir un effet.

\subsection{Les principaux effets}

\subsubsection{L'ombre portée}

L'\terme{ombre portée} ajoute une ombre derrière le texte. Elle est utile pour :

\begin{itemize}
    \item Détacher le texte du fond.
    \item Améliorer la lisibilité.
    \item Créer un effet de profondeur.
\end{itemize}

Paramètres :

\begin{itemize}
    \item \terme{Mode de fusion} : Multiplier par défaut.
    \item \terme{Couleur} : couleur de l'ombre (noir par défaut).
    \item \terme{Opacité} : transparence de l'ombre.
    \item \terme{Angle} : direction de la lumière.
    \item \terme{Distance} : distance de l'ombre.
    \item \terme{Extension} : extension de l'ombre.
    \item \terme{Taille} : flou de l'ombre.
\end{itemize}

\subsubsection{L'ombre interne}

L'\terme{ombre interne} ajoute une ombre à l'intérieur du texte. Elle est utile pour créer un effet de profondeur.

\subsubsection{Le contour}

Le \terme{contour} ajoute un contour autour du texte. Il est utile pour :

\begin{itemize}
    \item Détacher le texte du fond.
    \item Améliorer la lisibilité.
    \item Créer un effet de style.
\end{itemize}

Paramètres :

\begin{itemize}
    \item \terme{Taille} : épaisseur du contour.
    \item \terme{Position} : extérieur, intérieur, centré.
    \item \terme{Mode de fusion} : Normal par défaut.
    \item \terme{Couleur} : couleur du contour.
    \item \terme{Opacité} : transparence du contour.
\end{itemize}

\subsubsection{Le biseau et estampage}

Le \terme{biseau et estampage} crée un effet de relief sur le texte. Il est utile pour :

\begin{itemize}
    \item Créer un effet de volume.
    \item Créer un effet de métal, de pierre, etc.
\end{itemize}

\subsubsection{La lueur interne}

La \terme{lueur interne} ajoute une lueur à l'intérieur du texte. Utile pour créer un effet lumineux.

\subsubsection{La lueur externe}

La \terme{lueur externe} ajoute une lueur autour du texte. Utile pour créer un effet de halo, de néon.

\subsubsection{Le biseau et estampage (déjà vu)}

\subsubsection{La couleur, le dégradé et le motif}

Ces effets permettent de remplacer le remplissage du texte par :

\begin{itemize}
    \item Une \terme{couleur unie}.
    \item Un \terme{dégradé}.
    \item Un \terme{motif}.
\end{itemize}

\subsubsection{Le calque de réglage}

Le \terme{calque de réglage} permet d'appliquer des réglages (couleur, luminosité) au texte.

\subsection{Combiner les effets}

\subsubsection{Le principe}

On peut combiner plusieurs effets sur un même texte. Par exemple :

\begin{itemize}
    \item Ombre portée + contour.
    \item Biseau + dégradé.
    \item Lueur externe + contour.
\end{itemize}

\subsubsection{Précautions}

\begin{itemize}
    \item \terme{Ne pas abuser} des effets.
    \item \terme{Vérifier} la lisibilité.
    \item \terme{Hiérarchiser} les effets.
    \item \terme{Vérifier} le rendu final.
\end{itemize}

\subsection{Les bonnes pratiques}

\begin{itemize}
    \item \terme{Utiliser} les styles de calque pour le texte.
    \item \terme{Combiner} plusieurs effets avec parcimonie.
    \item \terme{Ajuster} les paramètres pour un résultat naturel.
    \item \terme{Vérifier} la lisibilité.
    \item \terme{Vérifier} le rendu sur différents fonds.
    \item \terme{Ne pas abuser} des effets (voir section 10).
    \item \terme{Enregistrer} les styles pour les réutiliser.
\end{itemize}

\subsection*{Conclusion de la section}

Les \terme{styles de calque} permettent d'appliquer des effets graphiques au texte (et à d'autres calques) de manière non destructive. On utilise l'ombre portée, le contour, le biseau, la lueur, le dégradé, le motif. On peut combiner plusieurs effets. Il est essentiel d'utiliser les styles de calque, de combiner avec parcimonie, d'ajuster les paramètres, de vérifier la lisibilité, de ne pas abuser des effets.

Maintenant que nous savons utiliser les styles de calque, nous allons, dans la section suivante, apprendre à \terme{créer des ombres portées}.

% ============================================================
\section{Créer des ombres portées}
% ============================================================

\subsection{L'importance des ombres portées}

L'\terme{ombre portée} est l'un des effets les plus utilisés pour le texte. Elle permet de :

\begin{itemize}
    \item \terme{Détacher} le texte du fond.
    \item \terme{Améliorer} la lisibilité.
    \item \terme{Créer} un effet de profondeur.
    \item \terme{Ajouter} du style.
    \item \terme{Simuler} une source de lumière.
\end{itemize}

Mais l'ombre portée peut aussi \terme{nuire} à la lisibilité si elle est mal utilisée (trop grande, trop foncée, mal orientée).

\subsection{Créer une ombre portée avec les styles de calque}

\subsubsection{Étapes}

\begin{enumerate}
    \item Sélectionner le calque de texte.
    \item Double-cliquer sur le calque (ou clic droit > \emph{Options de fusion}).
    \item Cocher \emph{Ombre portée}.
    \item Ajuster les paramètres :
    \begin{itemize}
        \item \terme{Mode de fusion} : Multiplier (par défaut) ou autre.
        \item \terme{Couleur} : noir, ou couleur plus sombre.
        \item \terme{Opacité} : 50\% à 75\% (par défaut).
        \item \terme{Angle} : direction de la lumière.
        \item \terme{Distance} : 2 à 10 px.
        \item \terme{Extension} : 0\% à 10\%.
        \item \terme{Taille} : 5 à 20 px.
    \end{itemize}
    \item Valider.
\end{enumerate}

\subsection{Ajuster les paramètres}

\subsubsection{Le mode de fusion}

Le \terme{mode de fusion} détermine comment l'ombre se mélange avec le fond :

\begin{itemize}
    \item \terme{Multiplier} : assombrit le fond (par défaut).
    \item \terme{Normal} : ombre opaque.
    \item \terme{Superposition} : augmente le contraste.
\end{itemize}

\subsubsection{La couleur}

La \terme{couleur} de l'ombre est souvent noire, mais on peut utiliser :

\begin{itemize}
    \item \terme{Une couleur plus sombre} que le fond.
    \item \terme{Une couleur complémentaire} au fond.
    \item \terme{Une couleur de la palette} de la composition.
\end{itemize}

\subsubsection{L'opacité}

L'\terme{opacité} détermine la transparence de l'ombre :

\begin{itemize}
    \item \terme{Opacité élevée} : ombre foncée.
    \item \terme{Opacité faible} : ombre légère.
\end{itemize}

\subsubsection{L'angle}

L'\terme{angle} détermine la direction de la lumière :

\begin{itemize}
    \item \terme{Angle cohérent} : toutes les ombres doivent avoir le même angle.
    \item \terme{Cohérence} : angle cohérent avec les autres éléments de la composition.
\end{itemize}

\subsubsection{La distance}

La \terme{distance} détermine la distance entre le texte et l'ombre :

\begin{itemize}
    \item \terme{Distance courte} : ombre proche.
    \item \terme{Distance longue} : ombre éloignée.
\end{itemize}

\subsubsection{La taille}

La \terme{taille} détermine le flou de l'ombre :

\begin{itemize}
    \item \terme{Taille faible} : ombre nette.
    \item \terme{Taille élevée} : ombre floue.
\end{itemize}

\subsection{Les ombres portées avancées}

\subsubsection{L'ombre portée avec contour}

\begin{enumerate}
    \item Ajouter une ombre portée.
    \item Ajouter un contour.
    \item Vérifier la cohérence.
\end{enumerate}

\subsubsection{L'ombre portée multiple}

\begin{enumerate}
    \item Ajouter une première ombre portée.
    \item Ajouter une seconde ombre portée (avec un angle différent).
    \item Vérifier la cohérence.
\end{enumerate}

\subsubsection{L'ombre portée avec dégradé}

\begin{enumerate}
    \item Ajouter une ombre portée.
    \item Ajouter un dégradé sur l'ombre (via les options avancées).
    \item Vérifier la cohérence.
\end{enumerate}

\subsection{Les bonnes pratiques}

\begin{itemize}
    \item \terme{Cohérence} : angle cohérent avec les autres ombres.
    \item \terme{Discrétion} : ombre légère, pas trop foncée.
    \item \terme{Lisibilité} : l'ombre ne doit pas nuire à la lisibilité.
    \item \terme{Non-destructivité} : utiliser les styles de calque.
    \item \terme{Vérifier} sur différents fonds.
    \item \terme{Ajuster} l'opacité et la taille.
    \item \terme{Ne pas abuser} : une seule ombre portée suffit souvent.
\end{itemize}

\subsection*{Conclusion de la section}

Les \terme{ombres portées} permettent de détacher le texte du fond, d'améliorer la lisibilité, de créer un effet de profondeur. On les crée avec les styles de calque (Ombre portée). On ajuste les paramètres : mode de fusion, couleur, opacité, angle, distance, taille. Il est essentiel de respecter la cohérence, la discrétion, la lisibilité, la non-destructivité.

Maintenant que nous savons créer des ombres portées, nous allons, dans la section suivante, apprendre à \terme{appliquer des contours et des dégradés}.

% ============================================================
\section{Appliquer des contours et des dégradés}
% ============================================================

\subsection{L'importance des contours et des dégradés}

Les \terme{contours} et les \terme{dégradés} sont deux effets graphiques très utilisés pour le texte :

\begin{itemize}
    \item Le \terme{contour} détache le texte du fond, améliore la lisibilité, ajoute du style.
    \item Le \terme{dégradé} ajoute de la couleur, de la profondeur, de la texture.
\end{itemize}

\subsection{Le contour}

\subsubsection{Créer un contour}

\begin{enumerate}
    \item Sélectionner le calque de texte.
    \item Double-cliquer sur le calque (ou clic droit > \emph{Options de fusion}).
    \item Cocher \emph{Contour}.
    \item Ajuster les paramètres :
    \begin{itemize}
        \item \terme{Taille} : 1 à 5 px.
        \item \terme{Position} : extérieur, intérieur, centré.
        \item \terme{Mode de fusion} : Normal.
        \item \terme{Couleur} : couleur du contour.
        \item \terme{Opacité} : transparence du contour.
    \end{itemize}
    \item Valider.
\end{enumerate}

\subsubsection{Le contour multiple}

On peut ajouter plusieurs contours sur un même texte :

\begin{enumerate}
    \item Ajouter un premier contour.
    \item Cliquer sur le bouton \emph{+} pour ajouter un second contour.
    \item Ajuster les paramètres du second contour (plus grand, autre couleur).
\end{enumerate}

\subsection{Le dégradé}

\subsubsection{Créer un dégradé}

\begin{enumerate}
    \item Sélectionner le calque de texte.
    \item Double-cliquer sur le calque (ou clic droit > \emph{Options de fusion}).
    \item Cocher \emph{Dégradé} (ou \emph{Incrustation dégradé}).
    \item Ajuster les paramètres :
    \begin{itemize}
        \item \terme{Mode de fusion} : Normal.
        \item \terme{Opacité} : 100\%.
        \item \terme{Dégradé} : cliquer sur le dégradé pour le modifier.
        \item \terme{Style} : linéaire, radial, angulaire, réfléchi, losange.
        \item \terme{Angle} : orientation du dégradé.
        \item \terme{Échelle} : taille du dégradé.
    \end{itemize}
    \item Valider.
\end{enumerate}

\subsubsection{Créer un dégradé personnalisé}

\begin{enumerate}
    \item Ouvrir le sélecteur de dégradé.
    \item Ajouter des arrêts de couleur.
    \item Définir les couleurs et les positions.
    \item Valider.
\end{enumerate}

\subsection{Combiner contour et dégradé}

\subsubsection{Le principe}

On peut combiner contour et dégradé sur un même texte :

\begin{itemize}
    \item Un \terme{remplissage dégradé} avec un \terme{contour} sombre.
    \item Un \terme{remplissage uni} avec un \terme{contour dégradé}.
    \item Un \terme{dégradé multiple} avec un \terme{contour multiple}.
\end{itemize}

\subsection{Les bonnes pratiques}

\begin{itemize}
    \item \terme{Utiliser} les styles de calque (non destructif).
    \item \terme{Ajuster} les paramètres avec parcimonie.
    \item \terme{Vérifier} la lisibilité.
    \item \terme{Respecter} la cohérence (couleurs, style).
    \item \terme{Combiner} contour et dégradé avec parcimonie.
    \item \terme{Vérifier} le rendu sur différents fonds.
    \item \terme{Ne pas abuser} des effets.
\end{itemize}

\subsection*{Conclusion de la section}

Les \terme{contours} et les \terme{dégradés} sont deux effets graphiques très utilisés pour le texte. On les crée avec les styles de calque. On ajuste les paramètres : taille, position, mode de fusion, couleur, opacité, style, angle, échelle. On peut combiner contour et dégradé. Il est essentiel d'utiliser les styles de calque, d'ajuster les paramètres, de vérifier la lisibilité, de respecter la cohérence, de ne pas abuser.

Maintenant que nous savons appliquer des contours et des dégradés, nous allons, dans la section suivante, apprendre à \terme{créer des effets de profondeur}.

% ============================================================
\section{Créer des effets de profondeur}
% ============================================================

\subsection{L'importance des effets de profondeur}

Les \terme{effets de profondeur} donnent du volume, du relief, de la dimension au texte. Ils sont utilisés pour :

\begin{itemize}
    \item \terme{Donner} du volume (biseau, ombre interne).
    \item \terme{Créer} une impression de 3D.
    \item \terme{Attirer} l'attention.
    \item \terme{Ajouter} du style.
\end{itemize}

\subsection{Les effets de profondeur}

\subsubsection{Le biseau et estampage}

Le \terme{biseau et estampage} crée un effet de relief sur le texte. Il simule une lumière venant d'un côté et une ombre de l'autre.

\begin{enumerate}
    \item Double-cliquer sur le calque de texte.
    \item Cocher \emph{Biseau et estampage}.
    \item Ajuster les paramètres :
    \begin{itemize}
        \item \terme{Style} : biseau interne, biseau externe, etc.
        \item \terme{Méthode} : lisse, sculpté, estompé.
        \item \terme{Profondeur} : intensité du relief.
        \item \terme{Taille} : taille du relief.
        \item \terme{Adoucir} : adoucissement du relief.
        \item \terme{Angle} : direction de la lumière.
        \item \terme{Altitude} : hauteur de la lumière.
        \item \terme{Lueur} : mode de fusion, couleur, opacité.
        \item \terme{Ombre} : mode de fusion, couleur, opacité.
    \end{itemize}
\end{enumerate}

\subsubsection{L'ombre interne}

L'\terme{ombre interne} ajoute une ombre à l'intérieur du texte. Elle renforce l'effet de profondeur.

\begin{enumerate}
    \item Double-cliquer sur le calque de texte.
    \item Cocher \emph{Ombre interne}.
    \item Ajuster les paramètres :
    \begin{itemize}
        \item \terme{Mode de fusion} : Multiplier.
        \item \terme{Couleur} : noir ou couleur sombre.
        \item \terme{Opacité} : 50\% à 75\%.
        \item \terme{Angle} : direction de la lumière.
        \item \terme{Distance} : 2 à 5 px.
        \item \terme{Extension} : 0\% à 10\%.
        \item \terme{Taille} : 5 à 10 px.
    \end{itemize}
\end{enumerate}

\subsubsection{Le biseau avec dégradé}

On peut combiner un biseau avec un dégradé pour un effet métallique, de pierre, etc.

\subsubsection{L'effet 3D}

Photoshop propose un effet 3D natif (\texttt{3D > Nouvel objet 3D à partir du calque}). Cet effet est plus complexe et moins utilisé aujourd'hui.

\subsection{Combiner les effets}

\subsubsection{Le principe}

On peut combiner plusieurs effets de profondeur :

\begin{itemize}
    \item \terme{Biseau + ombre interne}.
    \item \terme{Biseau + dégradé}.
    \item \terme{Biseau + ombre portée}.
\end{itemize}

\subsection{Les bonnes pratiques}

\begin{itemize}
    \item \terme{Utiliser} les styles de calque.
    \item \terme{Ajuster} les paramètres pour un résultat naturel.
    \item \terme{Vérifier} la lisibilité.
    \item \terme{Respecter} la cohérence (angle, couleur).
    \item \terme{Combiner} avec parcimonie.
    \item \terme{Vérifier} le rendu sur différents fonds.
    \item \terme{Ne pas abuser} des effets de profondeur.
\end{itemize}

\subsection*{Conclusion de la section}

Les \terme{effets de profondeur} donnent du volume, du relief, de la dimension au texte. On utilise le biseau et estampage, l'ombre interne, le biseau avec dégradé, l'effet 3D. On peut combiner plusieurs effets. Il est essentiel d'utiliser les styles de calque, d'ajuster les paramètres, de vérifier la lisibilité, de respecter la cohérence, de ne pas abuser.

Maintenant que nous savons créer des effets de profondeur, nous allons, dans la section suivante, apprendre à \terme{intégrer des textures}.

% ============================================================
\section{Intégrer des textures}
% ============================================================

\subsection{L'importance des textures}

Les \terme{textures} ajoutent de la richesse, de la matière, du réalisme au texte. Elles sont utilisées pour :

\begin{itemize}
    \item \terme{Donner} un aspect (métal, bois, pierre, tissu).
    \item \terme{Créer} un effet de matière.
    \item \terme{Ajouter} de la profondeur.
    \item \terme{Attirer} l'attention.
\end{itemize}

\subsection{Les méthodes d'intégration des textures}

\subsubsection{Le masque d'écrêtage}

\begin{enumerate}
    \item Placer une texture dans un calque au-dessus du texte.
    \item Clic droit sur le calque de texture > \texttt{Créer un masque d'écrêtage}.
    \item La texture n'apparaît que dans les lettres du texte.
    \item Ajuster la position, la taille, l'opacité de la texture.
\end{enumerate}

\subsubsection{Le mode de fusion}

\begin{enumerate}
    \item Placer une texture dans un calque au-dessus du texte.
    \item Changer le mode de fusion du calque de texture (par exemple, \emph{Superposition}, \emph{Écran}, \emph{Multiplier}).
    \item Ajuster l'opacité.
\end{enumerate}

\subsubsection{Le style de calque Motif}

\begin{enumerate}
    \item Double-cliquer sur le calque de texte.
    \item Cocher \emph{Incrustation de motif}.
    \item Choisir un motif dans la bibliothèque.
    \item Ajuster l'échelle, le mode de fusion, l'opacité.
\end{enumerate}

\subsubsection{Le style de calque Dégradé avec texture}

Le dégradé peut être combiné avec une texture pour un effet plus riche.

\subsection{Les types de textures}

\subsubsection{Les textures photographiques}

\begin{itemize}
    \item \terme{Métal} : acier, cuivre, or.
    \item \terme{Bois} : chêne, noyer, etc.
    \item \terme{Pierre} : marbre, granit.
    \item \terme{Tissu} : lin, soie, jean.
    \item \terme{Papier} : kraft, parchemin.
\end{itemize}

\subsubsection{Les textures graphiques}

\begin{itemize}
    \item \terme{Motifs} : rayures, pois, carreaux.
    \item \terme{Trames} : points, lignes.
    \item \terme{Grunge} : éraflures, taches.
\end{itemize}

\subsubsection{Les textures générées}

\begin{itemize}
    \item \terme{Bruit} : ajout de bruit pour un effet granuleux.
    \item \terme{Filtres} : nuages, fibres, etc.
\end{itemize}

\subsection{Les bonnes pratiques}

\begin{itemize}
    \item \terme{Choisir} une texture adaptée au message.
    \item \terme{Utiliser} un masque d'écrêtage pour intégrer la texture.
    \item \terme{Ajuster} le mode de fusion et l'opacité.
    \item \terme{Vérifier} la lisibilité du texte.
    \item \terme{Respecter} la cohérence de la composition.
    \item \terme{Ne pas abuser} des textures (surcharge).
    \item \terme{Vérifier} le rendu sur différents fonds.
\end{itemize}

\subsection*{Conclusion de la section}

Les \terme{textures} ajoutent de la richesse, de la matière, du réalisme au texte. On les intègre avec le masque d'écrêtage, le mode de fusion, le style de calque Motif. On utilise des textures photographiques, graphiques, générées. Il est essentiel de choisir une texture adaptée, d'utiliser un masque d'écrêtage, d'ajuster le mode de fusion et l'opacité, de vérifier la lisibilité, de ne pas abuser.

Maintenant que nous savons intégrer des textures, nous allons, dans la section suivante, apprendre à \terme{construire une composition publicitaire}.

% ============================================================
\section{Construire une composition publicitaire}
% ============================================================

\subsection{Les principes de la composition publicitaire}

Une \terme{composition publicitaire} combine texte, image, couleurs, effets pour transmettre un message et inciter à l'action. Elle repose sur :

\begin{itemize}
    \item Un \terme{message clair} : un seul message principal.
    \item Une \terme{hiérarchie visuelle} : le message principal est le plus visible.
    \item Une \terme{identité} : logo, couleurs, typographies cohérents.
    \item Un \terme{impact} : attractif, mémorable.
    \item Un \terme{appel à l'action} : clair, visible.
\end{itemize}

\subsection{Les éléments d'une composition publicitaire}

\subsubsection{Le visuel}

Le \terme{visuel} est l'image principale : photographie, illustration, graphique. Il doit être attractif, adapté au message, de qualité.

\subsubsection{Le titre}

Le \terme{titre} est le message principal. Il doit être :

\begin{itemize}
    \item \terme{Court} : 3 à 7 mots.
    \item \terme{Percutant} : fort, mémorable.
    \item \terme{Lisible} : grand, contrasté.
    \item \terme{Adapté} : au public et au message.
\end{itemize}

\subsubsection{Le sous-titre}

Le \terme{sous-titre} complète le titre. Il apporte une information supplémentaire.

\subsubsection{Le texte}

Le \terme{texte} donne les informations détaillées : description, avantages, mentions.

\subsubsection{Le logo}

Le \terme{logo} identifie la marque. Il doit être visible mais discret.

\subsubsection{L'appel à l'action}

L'\terme{appel à l'action} (ou CTA, Call To Action) indique au spectateur ce qu'il doit faire : acheter, s'inscrire, visiter, etc.

\subsection{Les étapes de la composition}

\subsubsection{Étape 1 : Définir le message}

Quel est le message principal ? Que doit retenir le spectateur ?

\subsubsection{Étape 2 : Définir le public}

À qui s'adresse la publicité ? Quel est son âge, son milieu, ses habitudes ?

\subsubsection{Étape 3 : Choisir le visuel}

Choisir une image attractive, adaptée au message et au public.

\subsubsection{Étape 4 : Définir la hiérarchie}

Définir l'ordre d'importance : titre, visuel, sous-titre, texte, logo, CTA.

\subsubsection{Étape 5 : Choisir la typographie}

Choisir 1 à 2 polices, créer une hiérarchie de tailles.

\subsubsection{Étape 6 : Choisir les couleurs}

Choisir une palette cohérente (2 à 4 couleurs).

\subsubsection{Étape 7 : Appliquer les effets}

Appliquer les effets graphiques avec parcimonie (ombres, contours, dégradés, textures).

\subsubsection{Étape 8 : Vérifier}

Vérifier la lisibilité, la cohérence, l'impact.

\subsection{Exemple de composition publicitaire}

\subsubsection{Publicité pour un parfum}

\begin{itemize}
    \item \terme{Visuel} : photo du flacon, sur fond sombre.
    \item \terme{Titre} : nom du parfum, en grand, en or.
    \item \terme{Sous-titre} : slogan, en petit, en blanc.
    \item \terme{Texte} : description, en petit, en gris.
    \item \terme{Logo} : en bas, en blanc.
    \item \terme{CTA} : « Découvrir », en bas, en or.
    \item \terme{Effets} : ombre portée sur le titre, dégradé doré.
\end{itemize}

\subsection{Les bonnes pratiques}

\begin{itemize}
    \item \terme{Un seul message principal}.
    \item \terme{Hiérarchie} claire.
    \item \terme{Typographie} limitée (1 à 2 polices).
    \item \terme{Couleurs} cohérentes (2 à 4).
    \item \terme{Effets} avec parcimonie.
    \item \terme{Vérifier} la lisibilité.
    \item \terme{Vérifier} l'impact.
    \item \terme{Tester} sur différents supports.
\end{itemize}

\subsection*{Conclusion de la section}

La \terme{composition publicitaire} combine texte, image, couleurs, effets pour transmettre un message et inciter à l'action. Elle repose sur un message clair, une hiérarchie visuelle, une identité, un impact, un appel à l'action. Les étapes sont : définir le message, définir le public, choisir le visuel, définir la hiérarchie, choisir la typographie, choisir les couleurs, appliquer les effets, vérifier. Il est essentiel d'avoir un seul message, une hiérarchie claire, une typographie limitée, des couleurs cohérentes, des effets avec parcimonie.

Maintenant que nous savons construire une composition publicitaire, nous allons, dans la dernière section de ce chapitre, apprendre à \terme{éviter la surcharge d'effets}.

% ============================================================
\section{Éviter la surcharge d'effets}
% ============================================================

\subsection{L'importance de la retenue}

Les effets graphiques sont tentants : on veut ajouter des ombres, des contours, des dégradés, des textures, des biseaux, des lueurs. Mais \terme{trop d'effets} nuit à la lisibilité, à la crédibilité, à l'élégance. Un texte surchargé d'effets est :

\begin{itemize}
    \item \terme{Illisible} : on ne peut plus lire le texte.
    \item \terme{Amateur} : l'effet est trop visible.
    \item \terme{Kitsch} : l'accumulation d'effets donne un résultat criard.
    \item \terme{Épuisant} : l'œil ne sait plus où regarder.
\end{itemize}

Un texte avec peu d'effets est souvent plus élégant, plus professionnel, plus efficace.

\subsection{Les principes de la retenue}

\subsubsection{Un effet par texte}

En général, un seul effet suffit. Par exemple :

\begin{itemize}
    \item Une ombre portée légère.
    \item Un contour fin.
    \item Un dégradé subtil.
\end{itemize}

\subsubsection{Des effets discrets}

Les effets doivent être \terme{discrets} :

\begin{itemize}
    \item \terme{Ombre portée} : opacité 30\% à 50\%, taille modérée.
    \item \terme{Contour} : taille 1 à 2 px.
    \item \terme{Dégradé} : subtil, pas trop contrasté.
    \item \terme{Texture} : légère, pas trop marquée.
\end{itemize}

\subsubsection{Des effets cohérents}

Les effets doivent être \terme{cohérents} avec l'ambiance :

\begin{itemize}
    \item \terme{Élégant} : effets discrets, sobres.
    \item \terme{Dynamique} : effets plus marqués, mais maîtrisés.
    \item \terme{Ludique} : effets colorés, mais cohérents.
\end{itemize}

\subsubsection{Des effets au service du message}

Les effets doivent \terme{servir le message}, pas l'inverse. Un effet qui nuit à la lisibilité est un effet raté.

\subsection{Les erreurs courantes}

\subsubsection{Accumulation d'effets}

Empiler ombre + contour + biseau + dégradé + texture donne un résultat surchargé.

\subsubsection{Effets excessifs}

Ombre trop foncée, contour trop épais, dégradé trop contrasté, biseau trop profond.

\subsubsection{Effets incohérents}

Effets qui ne correspondent pas à l'ambiance (par exemple, un biseau sur un texte élégant).

\subsubsection{Effets au détriment de la lisibilité}

Effets qui rendent le texte illisible (par exemple, ombre portée trop foncée, contour trop épais).

\subsection{Les bonnes pratiques}

\begin{itemize}
    \item \terme{Limiter} le nombre d'effets (1 à 2 par texte).
    \item \terme{Ajuster} les paramètres avec parcimonie.
    \item \terme{Vérifier} la lisibilité.
    \item \terme{Vérifier} la cohérence avec l'ambiance.
    \item \terme{Vérifier} le rendu sur différents fonds.
    \item \terme{Faire} une pause avant de finaliser.
    \item \terme{Demander} un avis extérieur.
    \item \terme{Simplifier} si nécessaire.
\end{itemize}

\subsection{Les cas particuliers}

\subsubsection{Les visuels intentionnellement surchargés}

Certains visuels sont intentionnellement surchargés (par exemple, certains styles rétro, certains visuels d'art). Dans ce cas, la surcharge est un choix artistique assumé.

\subsubsection{Les textes très petits}

Pour les textes très petits, les effets doivent être encore plus discrets, car ils nuisent à la lisibilité.

\subsubsection{Les textes sur fond chargé}

Pour les textes sur fond chargé, l'ombre portée ou le contour peuvent être nécessaires pour la lisibilité. Mais ils doivent rester discrets.

\subsection*{Conclusion de la section}

La \terme{surcharge d'effets} nuit à la lisibilité, à la crédibilité, à l'élégance. Il faut adopter une \terme{retenue} : un effet par texte, des effets discrets, des effets cohérents, des effets au service du message. Il est essentiel de limiter le nombre d'effets, d'ajuster les paramètres, de vérifier la lisibilité, de vérifier la cohérence, de faire une pause, de simplifier.

% ============================================================
\section*{Conclusion du chapitre}
\addcontentsline{toc}{section}{Conclusion du chapitre}
% ============================================================

Nous voici au terme du chapitre 12. Nous avons appris à maîtriser le texte et les effets graphiques dans Photoshop :

\begin{itemize}
    \item \terme{Créer du texte} : outils de texte (horizontal, vertical, masque), types (ponctuel, paragraphe, curviligne), vectorisation.
    \item \terme{Propriétés typographiques} : police, taille, style, interlignage, crénage, approche, couleur, alignement, styles.
    \item \terme{Hiérarchie de titres} : niveaux, moyens (taille, graisse, couleur, police, casse, position, espacement), construction.
    \item \terme{Styles de calque} : ombre portée, contour, biseau, lueur, dégradé, motif ; combinaison.
    \item \terme{Ombres portées} : création, paramètres (mode de fusion, couleur, opacité, angle, distance, taille).
    \item \terme{Contours et dégradés} : création, paramètres (taille, position, mode de fusion, couleur, opacité, style, angle, échelle).
    \item \terme{Effets de profondeur} : biseau, ombre interne, biseau avec dégradé, effet 3D.
    \item \terme{Textures} : masque d'écrêtage, mode de fusion, style de calque Motif ; types (photographique, graphique, généré).
    \item \terme{Composition publicitaire} : éléments (visuel, titre, sous-titre, texte, logo, CTA), étapes, bonnes pratiques.
    \item \terme{Surcharge d'effets} : erreurs, retenue, bonnes pratiques.
\end{itemize}

Le texte est l'un des éléments centraux du design graphique. Il porte le message, structure la composition, transmet une identité. Dans Photoshop, le texte est un calque spécifique, modifiable, auquel on peut appliquer des propriétés typographiques et des effets graphiques. Ces effets permettent de créer des visuels typographiques spectaculaires, utilisés dans la publicité, l'édition, le web, le packaging. Mais les effets sont à double tranchant : bien utilisés, ils subliment le texte ; mal utilisés, ils le rendent illisible, kitsch, amateur.

Ce chapitre t'a donné les bases du texte et des effets graphiques dans Photoshop. Mais la maîtrise vient avec la pratique. N'hésite pas à t'entraîner sur des projets variés : affiches, bannières, posts réseaux sociaux, covers. Chaque projet est un nouveau défi, et chaque défi est une occasion de progresser.

Dans le chapitre suivant, nous aborderons la \terme{création de supports publicitaires avec Photoshop}. Tu apprendras à concevoir une affiche, un flyer, un visuel pour les réseaux sociaux, une bannière, à intégrer des photographies et des textes, à organiser la hiérarchie visuelle, à décliner un visuel en plusieurs formats, à contrôler la qualité, à exporter les fichiers finaux. Ces compétences te permettront de créer des visuels professionnels pour la communication.

\vspace{2em}

\begin{center}
\textit{Fin du chapitre 12 --- Texte et effets graphiques dans Photoshop.}

\vspace{1em}

\textbf{Prochain chapitre : Chapitre 13 --- Création de supports publicitaires avec Photoshop.}
\end{center}